\section{Results}
\label{sec:results}

\subsection{Benchmark setup}
\curia{} was evaluated on a panel of 19 mammalian genomes, using human (hg38) as
the reference and the others as queries, spanning primates through the marsupial
outgroups: rhesus macaque (rheMac10), mouse (mm39), rat (rn7), rabbit (HLoryCun3),
cow (bosTau9), pig (susScr11), Bryde's whale (HLbalEde1), dromedary (HLcamDro2),
horse (equCab3), cat (felCat9), dog (canFam5), pangolin (HLmanJav2), flying fox
(HLpteVam2), hedgehog (eriEur2), armadillo (dasNov3), Asian elephant (HLeleMax1),
and three marsupials --- gray opossum (monDom5), Virginia opossum (HLdidVir1), and
tammar wallaby (HLnotEug3). Genome assemblies
were obtained from NCBI \citep{sayersDatabaseResourcesNational2022} and the UCSC
Genome Browser \citep{kentHumanGenomeBrowser2002}. Due to variability in
annotation quality across species, high-quality ncRNA annotations were used only
for human (GENCODE V48) and mouse (GENCODE VM37)
\citep{harrowGENCODEReferenceHuman2012}, while other species were analyzed without
relying on their native annotations. Human and mouse contain on the order of
${\sim}20{,}000$ annotated ncRNA transcripts, spanning both well-supported
elements (e.g., tRNAs, rRNAs, selected snoRNAs) and less-characterized transcripts,
particularly among lncRNAs. Evaluation was performed at both the locus level (for
short ncRNAs and lncRNA loci) and the substructure level (for island matches),
following the criteria defined in Section~\ref{sec:eval}.

\subsection{Overall performance}
We first evaluated agreement between predicted loci and existing annotations. In
the human-to-mouse evaluation (mouse being the query for which a comprehensive
reference annotation is available), among 2{,}994 predicted short ncRNA loci
52.4\% show any overlap with an annotated exon ($\geq 1$\,nt; any-overlap rate),
46.5\% meet a reciprocal-overlap recovery criterion (predicted locus covering
$\geq 50\%$ of the annotated transcript), and 26.0\% meet near-complete boundary
recovery ($\geq 99\%$ overlap). Among any-overlap predictions, the median overlap
is 98.9\% (mean 86.8\%), indicating that matched loci are typically well
localized and consistent with annotated boundaries.

These results suggest that matched loci are generally well localized when
annotation support is present. The substantial fraction of predictions without
annotation overlap is consistent with a combination of annotation incompleteness,
projection errors, and potential false positives. Notably, many non-overlapping
predictions retain strong syntenic support, suggesting that some may correspond to
unannotated or weakly annotated ncRNA loci.

\subsection{Short ncRNA performance and score behavior}
% RiNALMo human-mouse numbers (analysis/pair_numbers.py + compute_fig4_data.py).
We evaluated similarity between matched short ncRNA loci using MMD in embedding
space. For the human--mouse pair (2{,}994 loci), the MMD distribution is skewed
toward low values (median 0.142; Fig.~\ref{fig:mmd}A), with a substantial fraction
of near-zero values (32.5\% below 0.05), indicating consistent embedding
representations between matched reference and query loci.

Stratification by biotype reveals clear differences (Fig.~\ref{fig:mmd}C).
snoRNAs and tRNAs exhibit the lowest MMD values (median $\approx 0$), consistent
with their strong, well-characterized structural conservation
\citep{chanTRNAscanSE20Improved2021}. miRNAs and lncRNAs occupy an intermediate
range, whereas snRNAs, misc\_RNAs, and rRNA-derived loci show higher and broader
distributions, likely reflecting copy number, sequence divergence, and genomic
context.

MMD is strongly associated with annotation overlap. Predictions with low MMD
values ($\lesssim 0.1$) show high any-overlap rates with annotated loci
(${\sim}88\%$),
whereas higher MMD values show reduced overlap and increased dispersion.
Predictions with MMD $> 0.2$ rarely overlap annotations and are enriched for
ambiguous cases. Consistent with this, loci overlapping annotations are enriched
at low MMD values, while non-overlapping predictions are shifted toward higher
values (Fig.~\ref{fig:mmd}B), supporting the interpretation of MMD as a continuous
empirical measure of correspondence confidence.

Notably, a subset of high-confidence predictions (MMD $\approx 0$) does not
overlap existing annotations despite strong syntenic support, suggesting that
annotation incompleteness contributes to the observed gap.

To assess whether embedding-based similarity merely recapitulates sequence
conservation, we compared MMD to pairwise sequence identity for matched short
ncRNA loci. The two measures were strongly anticorrelated (Pearson $r = -0.80$;
Spearman $\rho = -0.85$; $n = 1{,}500$ sampled loci). However, this relationship is not
deterministic: at intermediate sequence identity ($\approx 40$--$60\%$), MMD
values span a broad range, indicating that embedding similarity is related to,
but not fully determined by, pairwise sequence identity in this regime.

\begin{figure}[t]
  \centering
  \includegraphics[width=\linewidth]{fig4_mmd.pdf}
  \caption{Behavior of embedding-based similarity (MMD) across short ncRNA loci
  for the human--mouse pair. \textbf{(A)} Sequence identity versus MMD for matched
  loci (strongly anticorrelated). \textbf{(B)} MMD distributions for loci that do
  versus do not overlap mouse annotation: annotation-overlapping loci concentrate
  at low MMD, non-overlapping loci shift higher. \textbf{(C)} MMD distribution by
  ncRNA biotype, ordered by median.}
  \label{fig:mmd}
\end{figure}

\subsection{Structural conservation analysis (``islands'')}
% Island match metric is cosine-Smith-Waterman distance d = 1/(1+score), not MMD.
We analyzed localized subregions (``islands'') within lncRNAs to assess
cross-species correspondence beyond the locus level. Islands are matched at
nucleotide resolution by Smith--Waterman alignment over the per-token cosine
similarity dotplot, and match quality is reported as a cosine-SW distance
$d = 1/(1+\text{score})$ (lower is better; Section~\ref{sec:islandalign}). Island
detection reveals that embedding-supported signal is sparse and localized. Most
lncRNAs contain a small number of discrete islands rather than uniform signal
across the full transcript (Fig.~\ref{fig:cores}A), indicating that the strongest
cross-species signal is concentrated in compact subregions. Across loci, islands
were detected in approximately 54\% of lncRNAs, with a median of 2 islands per
transcript.

To assess cross-species correspondence, we evaluated matches between reference and
query islands. At the gene level, 96.9\% of islanded lncRNAs exhibit at least one
cross-species match in at least one query species, indicating that detectable
correspondence is widespread but rarely spans entire transcripts. About half of
matched islands (50.1\%) overlap annotated ncRNA regions in the query genome, and
28.6\% show $\geq 50\%$ exonic overlap, suggesting that a substantial fraction of
candidate elements are not captured by current annotation resources.

Together, these observations support a modular pattern of cross-species
correspondence, with signal concentrated in discrete localized regions rather than
uniformly distributed across entire transcripts
\citep{hezroniPrinciplesLongNoncoding2015,rivasEvolutionaryConservationRNA2021,ulitskyEvolutionRescueUsing2016,mattickLongNoncodingRNAs2023,quinnUniqueFeaturesLong2016}.
This interpretation is broadly consistent with observations from well-characterized
lncRNAs such as NEAT1, XIST, and MALAT1, in which localized domains have been
linked to distinct architectural, spreading, or stability-related functions
\citep{wiluszTripleHelixStabilizes2012,yamazakiFunctionalDomainsNEAT12018,engreitzXistLncRNAExploits2013}.

\subsubsection{Multi-species conserved cores}
Here, ``conserved cores'' refers operationally to islands showing recurrent
high-confidence matches across multiple species within syntenically projected
loci. We observe a subset of islands that exhibit reproducible cross-species
matches (Fig.~\ref{fig:cores}B--C), with increasing divergence at greater
phylogenetic distance. Across all detected cores, 87\% are recovered in at least
two species, 67\% in five or more species, and approximately 2\% are consistently
detected across all 19 species, including the marsupial outgroups (opossums and
wallaby).

In contrast to locus-level overlap, which is often incomplete or ambiguous, these
cores exhibit stable cross-species correspondence at the level of localized
embedding similarity. To assess whether such matches could arise from selecting
the best of many candidate positions within a syntenic window, we compared each
matched query island against same-length windows tiled across its own projected
syntenic locus (a within-locus positional null; representative genome pairs). The
assigned position was the best-scoring one far more often than expected by chance,
and its advantage over the alternatives \emph{increased} with the number of
independent same-locus positions rather than collapsing toward chance, indicating
that recurrent cores are unlikely to be artifacts of a best-of-many maximum within
syntenic regions.

\subsection{Candidate conserved regions absent from current annotations}
After excluding regions overlapping or proximal to protein-coding genes, we
identify approximately 230 candidate intergenic loci potentially corresponding to
conserved ncRNA elements, exhibiting consistent embedding-space similarity across
$\geq 18$ of the 19 species (of 1{,}107 conserved-core lncRNA loci in total). The
full set of loci (all conservation categories) is provided as a supplementary
table in the repository:
\url{https://github.com/kirilenkobm/CURIA_pipeline/blob/master/paper/lncRNAs_with_conserved_cores.tsv}.
These loci
represent candidate elements and may include false positives despite filtering.
The persistence of these loci after filtering suggests that the observed signal is
unlikely to be primarily driven by coding-associated conservation, although
residual confounding cannot be excluded.

\subsection{Case studies}
To illustrate the behavior of \curia{}, we examined individual loci
(Fig.~\ref{fig:cases}). Projection of human SNORD57 to mouse identifies a nearly
exact match overlapping annotated Snord57, demonstrating precise localization
under strong structural conservation. Projection of human RNU6-7 intersects an
annotated mouse locus (Gm24859), indicating recovery of a corresponding genomic
region despite ambiguous or incomplete annotation. Projection of human vault RNA
(ENSG00000199990) identifies a corresponding locus overlapping mouse Vaultrc5; the
matched region differs in length from the human transcript, yet \curia{} localizes
the most similar segment within the syntenic region.

\begin{figure}[t]
  \centering
  \includegraphics[width=\linewidth]{fig5_cases.png}
  \caption{Recovery of short ncRNA loci across species using \curia{}. Genome
  browser snapshots (UCSC) illustrating representative examples of short ncRNA
  projection from human (hg38) to mouse (mm39): SNORD57 (exact match), RNU6-7
  (recovery under incomplete annotation), and vault RNA (length divergence,
  partial correspondence).}
  \label{fig:cases}
\end{figure}

\begin{figure}[t]
  \centering
  \includegraphics[width=\linewidth]{fig6_cores.pdf}
  \caption{Modular and heterogeneous cross-species correspondence of lncRNA cores.
  \textbf{(A)} Cross-species matching of core regions for representative lncRNAs
  (NEAT1, MIAT, MALAT1, XIST, NORAD); rows are species, columns are reference cores ordered along
  the transcript, color indicates match distance (cosine-SW distance $d=1/(1+s)$). \textbf{(B)} Distribution of
  core reproducibility across species. \textbf{(C)} Median island-level match
  distance across species, increasing with phylogenetic distance.}
  \label{fig:cores}
\end{figure}

\subsection{Error analysis}
The dominant failure modes arise from limitations in projection, annotation, and
signal strength. Chain-based projection defines the search space and acts as a
hard constraint: if the corresponding region is not captured by genome alignments,
it cannot be recovered. Incomplete and imprecise reference annotations further
limit recoverable signal, particularly for lncRNAs. A substantial fraction of loci
exhibit weak or diffuse embedding signal, producing no detectable islands or only
fragmented matches. Overall, most observed failures are associated with missing or
ambiguous signal and projection constraints, rather than instability of the
downstream matching procedure.

\subsection{Runtime and scalability}
% TODO(authors): confirm the exact GPU/CPU used for the reported panel (runs
% spanned more than one machine). Wall-time range below is measured from the 19
% completed run logs: min 1h09m, median ~2h27m, max 4h18m.
We evaluated \curia{} on the 19 genome-pair panel. On a single GPU, a full
genome-pair run completes in roughly one to four hours (median ${\sim}2.5$\,h
across the 19 pairs), scaling approximately linearly with annotation size.
Embedding inference is performed on GPU with batched execution, while downstream
matching and alignment are CPU-bound; both reference and query island detection
are restricted to syntenically projected regions, substantially reducing the
effective search space.
